\section{Introduction}
\label{sec:intro}

The preceding papers in this series constructed a Machine of State
that enforces Montesquieu's Separation of Powers in software.
Paper~1~\cite{btv2026} constrained the Executive with linear types;
Paper~2~\cite{btv-t-repo} created an immutable evidentiary memory;
Paper~3~\cite{btv-ar2026} enabled privacy-preserving auditing;
Paper~4~\cite{btv-econ} justified the economics of compliance; and
Paper~5~\cite{btv-cc} unified these into a theory of Constitutional
Code. The result is an institution whose constitution is enforced not
by social contract but by the mathematical properties of its execution
substrate.

There is a flaw in this edifice: \textbf{time}.

A perfect enforcement of today's law can become the perfect
perpetuator of tomorrow's injustice. If the type system is immutable,
the institution cannot adapt to new ethical norms, new scientific
understanding of algorithmic harm, or new democratic mandates. If the
type system is mutable by a simple regulatory vote, it is vulnerable
to \emph{regulatory capture}~\cite{stigler1971}---the very pathology
it was designed to prevent. We call this the \emph{Capture Paradox}:

\begin{itemize}
  \item To prevent capture, the constitution must be rigid
  (immutable types).
  \item To prevent obsolescence, the constitution must be flexible
  (upgradeable types).
  \item A constitution that is both rigid and flexible requires a
  \emph{formal process of change} that is itself capture-resistant.
\end{itemize}

This paper resolves the Capture Paradox by formalising the
\emph{process of constitutional change} as a first-class object in
the type system. The key insight, drawn from constitutional law, is
that not all rules are equal: some rules define the game, and some
rules define who gets to change the rules. The former we call
\emph{Stone Clauses}; the latter we formalise as the
\emph{Constitutional Amendment Protocol}.

\paragraph{Contributions.}
\begin{enumerate}
  \item \textbf{Temporal Type Versioning
  (\S\ref{subsec:versioning}):} a formal mechanism for versioning
  constitutional types so that historical decisions are always
  verifiable under the law that governed them at the time of
  execution, implementing the principle of \emph{nullum crimen sine
  lege} in software.
  \item \textbf{Stone Clauses and Tripartite Ratification
  (\S\ref{subsec:threshold}):} a classification of amendments into
  ordinary legislation (single-branch authorisation) and
  constitutional amendments (all three branches required), with a
  formal proof of capture-resistance.
  \item \textbf{Sunset Clauses as Linear Resources
  (\S\ref{subsec:sunset}):} encoding legislative authority as
  a finite-lifetime linear token, forcing periodic democratic
  renewal and making \emph{Constitutional Interregnum} a structural
  property of the system.
  \item \textbf{Security Analysis (\S\ref{sec:security}):}
  formal characterisation of two primary attack vectors---Hostile
  Takeover and Zombie Institution---and proofs of resistance
  derived from Amendment Soundness and Mandatory Renewal.
\end{enumerate}

Together, these contributions close the gap between static
constitutional enforcement (Papers 1--5) and dynamic institutional
evolution, completing the framework for Verifiable Institutions.
